\section{Course Schedule II}
\label{sec:graph-course-schedule-2}

\problemstatement{Given $\textit{numCourses}$ and prerequisite pairs
$[a,b]$ as in Section~\ref{sec:graph-course-schedule-1}, return
\textbf{one valid order} to take all courses, or an empty array if no
valid order exists.}

\begin{patternbox}
\emph{``Return a valid order respecting prerequisites''} is a direct
restatement of topological sort
(Sections~\ref{sec:graph-topo-sort-dfs}--\ref{sec:graph-topo-sort-kahn}):
build the same prerequisite graph as Course Schedule I, then run either
topological sort algorithm. Kahn's algorithm is the natural choice here
since it detects the no-valid-order case (a cycle) \textbf{for free},
via the same processed-count check from Section~\ref{sec:graph-cycle-directed}.
\end{patternbox}

\subsection*{Reduction}

Build the graph identically to Course Schedule I ($b \to a$ for every
pair $[a,b]$), run Kahn's algorithm, and return the resulting order --
but only if every node was processed; otherwise return an empty array
(a cycle means no valid order exists).

\subsection*{C++ implementation}

\begin{lstlisting}
#include <bits/stdc++.h>
using namespace std;

vector<int> findOrder(int numCourses, vector<vector<int>>& prerequisites) {
    vector<vector<int>> adj(numCourses);
    vector<int> inDegree(numCourses, 0);

    for (auto& p : prerequisites) {
        int a = p[0], b = p[1];
        adj[b].push_back(a);
        inDegree[a]++;
    }

    queue<int> q;
    for (int i = 0; i < numCourses; i++) {
        if (inDegree[i] == 0) q.push(i);
    }

    vector<int> order;
    while (!q.empty()) {
        int node = q.front();
        q.pop();
        order.push_back(node);

        for (int neighbor : adj[node]) {
            inDegree[neighbor]--;
            if (inDegree[neighbor] == 0) q.push(neighbor);
        }
    }

    if ((int)order.size() != numCourses) return {}; // cycle: no valid order
    return order;
}

int main() {
    vector<vector<int>> prerequisites = {{1, 0}, {2, 0}, {3, 1}, {3, 2}};
    vector<int> order = findOrder(4, prerequisites);
    for (int x : order) cout << x << " ";
    cout << endl; // 0 1 2 3  (one valid order)
    return 0;
}
\end{lstlisting}

\subsection*{Dry run}

\dryrun{Take $\textit{numCourses}=4$,
$\textit{prerequisites}=[[1,0],[2,0],[3,1],[3,2]]$ (course $3$ needs
both $1$ and $2$; both $1$ and $2$ need $0$).}

Edges: $0\to1,0\to2,1\to3,2\to3$. In-degrees: $0$:$0$, $1$:$1$,
$2$:$1$, $3$:$2$. Initial queue: $[0]$. Pop $0$: order $[0]$; decrement
$1\to0$ (push), $2\to0$ (push). Queue: $[1,2]$. Pop $1$: order
$[0,1]$; decrement $3\to1$ (not $0$ yet). Pop $2$: order $[0,1,2]$;
decrement $3\to0$ (push). Pop $3$: order $[0,1,2,3]$. Size $4 =
\textit{numCourses}$: valid order found.

\begin{complexitybox}
\textbf{Time Complexity.} $O(V+E)$, inherited unchanged from Kahn's
algorithm.
\textbf{Space Complexity.} $O(V+E)$ for the adjacency list and
in-degree array.
\end{complexitybox}

\begin{takeawaybox}
Course Schedule I and II together demonstrate the cleanest possible
illustration of this chapter's two core tools applied side by side:
the exact same prerequisite graph answers ``is it possible'' via cycle
detection and ``what is an order'' via topological sort, with no
additional derivation needed for either once the graph is built.
\end{takeawaybox}
